\section{Construction of the cross-country sectoral exposure database}
\label{app:database_construction}

This section implements the design described in Section~\ref{app:database-design}. The construction starts from published sector and instrument controls, allocates the detail needed for bilateral transmission, and records every reconciliation. Amounts are EUR billions unless stated otherwise. The final matrix, intermediate controls and long exposure register accompany the replication files.

\subsection{Target object, coverage and reference date}

Let $\mathcal C=\{\mathrm{DEU},\mathrm{FRA},\mathrm{ITA},\mathrm{ESP}\}$. Each country has five sectors: non-financial corporations (S11), banks (S122), nonbank financial institutions (S12nbfi), general government (S13), and households including non-profit institutions serving households (NPISHs) where the source combines them (S14). Nonbank financial institutions comprise S123--S129, including money market funds (MMFs). Banks exclude MMFs and central banks.

The reporting matrix $\mathbf W$ has creditor rows and debtor columns. It contains these twenty country--sector nodes, an aggregate central-bank node EA:S121, the rest of the world ROW:S2, and two classification accounts, UNALLOC:ISSUER and UNALLOC:HOLDER. It is therefore a $24\times24$ matrix. The classification accounts retain positions with an unidentified sector or holder; they are not additional economic sectors. Country information that survives in the source remains in the exposure register even when the matrix uses an unallocated account.

For model exposition and the full sector-network view, ROW:S2, UNALLOC:ISSUER and UNALLOC:HOLDER are grouped into a single passive node, Other counterparties. The resulting matrix has 22 accounts and retains the Eurosystem separately. Aggregation sums gross creditor and debtor positions without netting them. It preserves each identified bilateral claim, total bank assets and liabilities represented in the matrix, and the instrument-specific distinction between ordinary bank recovery and fixed settlements. The reporting matrix and exposure register retain the geographical and classification detail, since an unidentified domestic holder is not a foreign investor.

The retained instruments are bank loans, deposits and debt securities, and nonbank debt-security holdings and deposits with banks. Government loans remain in the separate balance-sheet reconciliation because the sovereign shock concerns debt securities. Equity, fund shares, derivatives and other assets are not reconstructed as a complete network of promised debt payments. Statistical remaining liabilities are reported separately by instrument in the register, including when they enter an unallocated matrix cell.

The database is a snapshot at 31 December 2024. Quarter-start dates such as \texttt{2024-10-01} label end-quarter positions. December monthly stocks and fourth-quarter stocks are alternatives, never additive observations. The annual Portfolio Investment Positions (PIP) observations used here refer to the same year end. No earlier quarterly position is inferred from the annual observation for this construction. ECB and IMF source snapshots were retrieved on 21 October 2026; Eurostat inputs retain the recorded extraction vintage. The replication files identify the source vintage for each input.

\subsection{Source hierarchy and construction sequence}

Table~\ref{tab:data_source_summary} summarizes the source hierarchy \citep{ECBData,EurostatData,IMFPIP}. BSI denotes Balance Sheet Items; SHSS denotes Securities Holdings Statistics; PIP was formerly the Coordinated Portfolio Investment Survey (CPIS). Source-table frequencies M and Q mean monthly and quarterly; financial-account codes F2, F3 and F4 denote currency and deposits, debt securities and loans. BIS consolidated claims and nationality-based domestic exposures do not enter the matrix. Worldwide group exposures cannot be spliced into resident-office balance sheets without a perimeter adjustment. They may serve as benchmarks, with their consolidation and counterparty definitions stated separately.

\begin{table}[H]
\centering
\floatcaption{Sources and their roles in the residence-based construction}{}
\label{tab:data_source_summary}

\small
\renewcommand{\arraystretch}{1.12}
\begin{tabularx}{\textwidth}{@{}>{\raggedright\arraybackslash}p{2.45cm}>{\raggedright\arraybackslash}p{2.9cm}>{\raggedright\arraybackslash}X@{}}
\toprule
Source & Instruments and frequency & Retained role \\
\midrule
ECB BSI & Loans, deposits, securities and balance-sheet totals; M/Q & Resident S122 controls; borrower and depositor sectors; bank and central-bank counterparties; related-office margins. \\
ECB SHSS & Debt securities; Q & Domestic bank and MMF holdings, with face and market values distinguished. \\
Eurostat financial accounts & Financial stocks; Q & Domestic nonbank holdings, government and central-bank portfolios, and domestic/nonresident issuer margins. \\
IMF PIP/CPIS & Portfolio debt securities; annual observation & Bilateral foreign portfolios and reported holder/issuer sector margins; constrained completion of missing sector detail. \\
Eurostat government debt & Consolidated nominal F3 stock; Q & Sovereign debt-security totals. Loans and other public-debt instruments are outside this anchor. \\
ECB ILM & Monetary-policy lending; M & Internal Liquidity Management statistics: refinancing diagnostic, not an addition to central-bank funding already recorded from BSI. \\
\bottomrule
\end{tabularx}
\end{table}

Construction proceeds in five steps: (i) establish resident bank instrument totals; (ii) allocate loan, deposit and securities counterparties; (iii) reconcile interbank positions and remove identified related-office loans; (iv) reconcile securities to issuer stocks while retaining unknown holders separately; and (v) reconcile selected claims to complete bank balance sheets. The source layers and all intermediate adjustments are retained. Aggregate sectors such as S1 discipline allocations but are never added to their component sectors as network nodes.

\subsection{Bank loans and deposit funding}

We first establish what banks lend and who funds them, using the same resident-bank controls on both sides. Three distinctions matter: bank versus money-market-fund reporting, the sector of private borrowers, and the sector of depositors. The adjustments below allocate those controls; they do not add a second layer of bank assets or funding.

\paragraph{Separate banks from money market funds.}
BSI reporting sector B identifies S122; sector A comprises monetary financial institutions (MFIs) excluding the European System of Central Banks (ESCB) and includes MMFs; sector F identifies MMFs. Where only sector-A bilateral detail is public, it is allocated within a matching sector-B instrument and counterparty class. For reporting country $i$, counterparty country $j$, instrument $k$ and counterparty class $g$, let $T_i^{k,g,\mathrm{S122}}$ and $T_i^{k,g,\mathrm{MFI}}$ be the corresponding published portfolio totals, and $w_{ij}^{k,g,\mathrm{MFI}}$ the bilateral MFI position. Define
\begin{equation}
 \rho_i^{k,g}=\frac{T_i^{k,g,\mathrm{S122}}}{T_i^{k,g,\mathrm{MFI}}},
 \qquad \widehat w_{ij}^{k,g}=\rho_i^{k,g}w_{ij}^{k,g,\mathrm{MFI}}.
 \label{eq:mmf-separation}
\end{equation}
This is an allocation of an observed total, not an observation of the bilateral S122 position. 

The register retains the MFI value, the allocation method and available bounds. For example, France's S122/MFI private-loan ratio is about 0.996: treating the MFI detail as a bank observation would retain some MMF lending. 

For outside-euro-area loans and securities, published MMF instrument totals are subtracted directly when available.

\paragraph{Identify loan borrowers.}
Domestic and bilateral sample-country loan observations distinguish nonfinancial corporations (NFCs) and households/NPISHs. The financial borrower residual is calculated within the private-sector loan class, after excluding government loans. Missing bilateral government loans are allocated using the observed foreign-euro-area government share within the nonbank envelope. Unassigned euro-area borrowers are retained in UNALLOC:ISSUER. Household loans use the reported household borrower-sector observations.

\paragraph{Identify deposit funding.}
Household and NFC deposits are taken from their reported depositor-sector cells, with the bank/MMF adjustment in \eqref{eq:mmf-separation} where necessary. Foreign financial and government depositors use the available disjoint sector breakdowns. Missing domestic financial/government detail is recovered from euro-area sector totals less reported foreign-euro-area positions. Any unclassified foreign depositor budget is allocated within the observed domestic residual, with the adjustment and bounds recorded. 

Domestic MMF depositors are separated from other MFIs and assigned to S12nbfi.

Outside-euro-area deposit funding is allocated between S122 and MMFs using their external-liability totals where the deposit split itself is unpublished. The remaining known external liabilities stay in an outside funding component. Unassigned euro-area deposits and other liabilities remain explicit. Household funding comprises measured deposits and separately reconciled bank securities.

\subsection{Interbank loans: reconcile mirrors, then remove related offices}

A loan from one bank is a liability of another. Asset-side reports and depositor-side mirrors provide complementary information on that same claim, but their public breakdowns need not agree cell by cell. We anchor the construction in loan asset totals and use funding mirrors and bilateral envelopes to allocate them. Identified related-office positions are removed only after the gross network has been formed.

The interbank loan construction covers all twenty countries belonging to the euro area at end-2024, before aggregating nonsample countries into ROW. Otherwise, exposures to the rest of the euro area would be forced into the four-country block. Central-bank counterparties are excluded from bank--bank loans and retained in EA:S121.

Let $N_{ij}$ denote gross loans from banks in country $i$ to banks in country $j$. The source provides bank loan asset margins, deposit-liability mirrors and, for many country pairs, broader bilateral MFI envelopes. Bank asset margins determine row totals. Foreign deposit mirrors are first rescaled to the same aggregate total. Available bilateral MFI values provide upper bounds; missing links receive a borrower-size prior based on foreign bank deposit margins. Domestic bank loan positions supply the diagonal, with the reporting-bank adjustment in \eqref{eq:mmf-separation}.

The allocation is a nonnegative, bounded weighted projection. Let $\mathbf{P}$ be the prior matrix, $\mathbf{r}$ the vector of row margins, $\mathbf{c}$ the vector of column targets and $\mathbf{U}$ the matrix of cell caps. For allocation matrix $\mathbf{Z}$, a small positive weight floor $\varepsilon$, a vector of ones $\mathbf1$ and transpose denoted by $\top$, the basic problem is
\begin{equation}
 \min_{\mathbf{Z}\geq0}\ \sum_{ij}\frac{(Z_{ij}-P_{ij})^2}{\max(P_{ij},\varepsilon)}
 \quad\text{s.t.}\quad
 \mathbf{Z}\mathbf1=\mathbf{r},\quad \mathbf{Z}^\top\mathbf1=\mathbf{c},\quad Z_{ij}\leq U_{ij}.
 \label{eq:bounded-allocation}
\end{equation}
When mirror columns are incompatible with the row margins and caps, rows and caps remain hard constraints and column deviations enter a weighted least-squares penalty. The register reports the resulting asset/liability mirror discrepancies.\footnote{Amounts are in EUR bn. The implementation uses $\varepsilon=10^{-5}$. With soft columns, the objective is $10^{-3}$ times the prior term in \eqref{eq:bounded-allocation} plus $\sum_j[(\mathbf{Z}^\top\mathbf1)_j-c_j]^2/\max(c_j,10^{-3})$. The positive denominator floor keeps the objective weights finite. All fitted margins and discrepancies are exported.}

BSI items A20G and L20G identify related-office loan/deposit margins. A related-office matrix $\mathbf{G}$ is fitted within the gross matrix, $0\leq G_{ij}\leq N_{ij}$. Identified foreign asset-side margins are removed in full; funding mirrors are reconciled within the gross bounds. The domestic removal uses the common observed subset, bounded by the adjusted related-office asset margin, the funding margin and the gross domestic position. Net interbank loans are
\[
 L_{ij}^{BB}=N_{ij}-G_{ij}.
\]
Missing auxiliary-country margins are estimated from available MFI/central-bank detail or median observed euro-area counterparty/related-office shares and flagged separately. Related securities and outside-euro-area related-office positions without a public split are not assumed absent.

Remaining domestic claims between independent banks are retained. Residence-based removal of identified related-office positions does not constitute worldwide group consolidation or full national-sector consolidation. The replication outputs also provide gross-related-office and zero-domestic-diagonal accounting variants. The model retains the adjusted gross domestic claims; a separate clearing run cancelling domestic bank claims and equal liabilities removes multiplicity at the reference parameters (Section~\ref{app:residence-model-mapping}).

\subsection{Bank securities: portfolio controls and issuer allocation}

The bank securities block follows a source hierarchy within observed portfolio totals. Sovereign holdings are treated first because they determine direct exposure to the model's initiating shock. Bank-issued and private securities are then allocated within their own controls.
S122 securities totals distinguish MFI, government and private-sector issuers. Domestic SHSS face-value holdings are preferred for sovereign exposures where public (Germany, France and Spain). Italy's domestic sovereign holding is allocated from MFI geography within its S122 government-security control. Foreign government portfolios use residence-based CPIS sector priors where available and bilateral MFI nonbank-security positions elsewhere, subject to the latter envelopes. Private portfolios use the corresponding residual nonbank-security envelopes; NFC/NBFI detail uses domestic SHSS or CPIS/issuer-stock priors. Allocations to nonsample euro-area issuers enter ROW. An allocation that cannot be reconciled to an issuer remains unallocated while the bank's instrument total is preserved.

BSI and SHSS bank-security positions are not interchangeable. BSI excludes retained own securities from both holdings and issuance, whereas SHSS holdings can include them. Domestic bank-to-bank securities therefore use BSI-matched class and geography; SHSS is retained as a benchmark. Differences can also reflect valuation and allocation, so the full gap is not identified as own retention \citep{ECBBSI2024}.

\subsection{Domestic securities held by nonbank sectors}
Domestic F3 holdings use Eurostat quarterly stock assets: frequency Q, \texttt{stk\_flow=STK}, \texttt{finpos=ASS}, and \texttt{unit=MIO\_EUR}. Transactions, other changes and liability observations are excluded. Household/NPISH sectors are mapped jointly where published that way. For financial subsectors, a reported S125--S127 aggregate takes precedence over its components; both are never added.

The domestic Eurostat holder slice contains no public S123 F3 observations. Domestic MMF holdings are therefore added from SHSS, with missing sector detail allocated within the domestic SHSS aggregate and BSI MMF portfolio controls. They belong to S12nbfi. No MMF top-up is added to foreign CPIS financial portfolios, which already include MMFs. Government and available central-bank domestic portfolios are retained from financial accounts. Where a central-bank sovereign holding is unpublished, the MFI holder aggregate less SHSS bank and MMF market holdings supplies a flagged residual when the three inputs are available.

\subsection{Foreign securities held by nonbank sectors}

Domestic holdings are supplemented with foreign securities positions from the IMF's residence-based portfolio survey. Its country-pair totals are more complete than its joint holder--issuer sector cells. We retain published compatible cells and complete the missing sector detail within those totals.

PIP debt-security assets (indicator \texttt{P\_F3\_P\_USD}) are classified by holder residence and sector and issuer residence and sector. The financial aggregate S12R includes insurance corporations and pension funds, MMFs and other financial institutions. We use S12R where available and its disjoint components only as alternatives, independently on both axes \citep{IMFCPIS2017}. This avoids adding a parent sector to its children.

Year-end USD positions are converted at the closing ECB exchange rate: 1.0389 USD per euro on 31 December 2024. Quarterly-average FX is not used for this stock conversion. Source observations remain recoverable in their original units.

For a country pair, let $T$ be total-economy holdings, $r_h$ the available holder-sector totals and $c_s$ the issuer-sector totals. The uncovered holder amount $T-\sum_h r_h$ enters an unallocated-holder row. Missing issuer sectors receive the residual to $T$, with relative weights from their nonresident debt-security liabilities. An observed zero issuer stock therefore does not receive a positive allocation merely to complete the table.

Reported joint holder--issuer cells are fixed when compatible with these margins. The remaining cells are allocated by \eqref{eq:bounded-allocation}, starting from the independence prior $r_hc_s/T$, after subtracting the fixed cells from each margin. In the present vintage, no proportional shrinking of reported holder margins or adjustment of reported joint cells is required at this stage. Incompatible margins or fixed cells would trigger a separately flagged reconciliation. This is a completion of missing sector detail, not an assertion that PIP publishes a full bilateral sector table.

The completed nonbank portfolios enter the matrix. Bank-holder portfolios serve as allocation priors only: BSI-controlled bank rows replace them, so they are never added as a second bank-security layer. Subsequent issuer reconciliation can adjust retained holdings. Accordingly, PIP country-pair totals are preserved in the intermediate portfolio table, not imposed on the final mixed-source matrix.

\subsection{Final checks against issuer stocks and outside positions}
\label{app:row}

The final step checks the holdings assembled from different sources against the securities outstanding. It aligns domestic and foreign valuation margins, resolves overfilled issuer columns and records unidentified holders explicitly. Bank asset totals and issuer totals are both controls; a discrepancy must be reconciled, not hidden in a negative outside position.

\paragraph{Domestic and foreign issuer margins.}
For each issuer, indexed by $u$, Eurostat counterpart S2 is the nonresident liability component, not the total amount outstanding. Domestic holdings are recovered from the assets of all resident holders (holder sector S1) against that issuer sector. These distinct margins discipline domestic and foreign portfolios.

Government debt securities use the consolidated nominal debt-security control from the Excessive Deficit Procedure (EDP). Holdings of domestic S13 bonds by S13 itself are removed before forming the domestic margin, consistently with consolidation. Bank-issued securities use the S122 BSI issuance control; other issuers use their financial-account security stocks. Let $M_u^d$ and $M_u^f$ be domestic and foreign reference stocks and $Q_u$ the issuer control. The common conversion factor is
\begin{equation}
 \nu_u=\frac{Q_u}{M_u^d+M_u^f},
 \qquad Q_u^d=\nu_u M_u^d,\qquad Q_u^f=\nu_u M_u^f.
 \label{eq:issuer-valuation}
\end{equation}
National-account and CPIS market holdings are converted with this issuer-average factor where needed. Existing ECB face/book inputs retain their source basis. Thus the nominal total is exact, but some individual bank cells remain statistical/book proxies rather than independently verified face amounts. A source-level valuation flag and an accounting-to-promised-payment bridge remain necessary for model use.

If represented holdings exceed a geographic issuer control, imputed positions are reduced first. If that is insufficient, the remaining holdings are reduced proportionally and relabelled reconciled. Raw values and adjustments stay in the register. Unmatched bank security assets return to UNALLOC:ISSUER, preserving the bank instrument control. Thus matching issuer totals can alter individual holdings, and the final matrix distinguishes reconciled cells from source observations.

\paragraph{ROW and unallocated accounts.}
For issuer $u$, let $\mathcal H_u^d$ and $\mathcal H_u^f$ contain represented domestic and foreign holders, respectively, and let $W_{vu}$ be holder $v$'s position in the reporting matrix. With $\mathrm{U_H}$ denoting the unallocated-holder account and ROW the rest of the world, the residuals are
\[
 W_{\mathrm{U_H},u}=Q_u^d-\sum_{v\in\mathcal H_u^d}W_{vu},
 \qquad
 W_{\mathrm{ROW},u}=Q_u^f-\sum_{v\in\mathcal H_u^f}W_{vu},
\]
after the relevant caps. The first residual is an unidentified domestic holder; the second is outside-sample foreign ownership. Separately, a CPIS holder with known country but unknown sector retains that country in the register. The two types of missing classification are not interpreted as household wealth.

ROW as debtor aggregates observed or allocated bank claims on nonsample euro-area countries and non-euro-area counterparties. The selected nonbank rows cover domestic and sample-country foreign securities; they are not closed to comprehensive world asset portfolios. Their empty ROW-debtor cells do not establish zero nonbank foreign wealth. Likewise, the central-bank node combines measured euro-area central-bank funding/deposit links and available domestic central-bank portfolios in sample-country securities; it is not a complete consolidated Eurosystem balance sheet. Monetary-policy refinancing from ILM is a diagnostic subset, not an extra creditor claim added to the BSI central-bank counterparty position. No capital-key allocation of an aggregate monetary-policy securities portfolio is used in this vintage.

\subsection{Register, missing values and reproducibility}

Each position records creditor and debtor, instrument, raw/reference and retained amounts, source keys, method, available bounds, inclusion in the matrix and one of five statuses: observed, derived, imputed, reconciled or unallocated. A reference amount can be a portfolio control rather than an observed bilateral position. Source files, metadata, retrieval queries and hashes are retained separately from the constructed exposures.

An unpublished observation is not replaced by an observed zero. Published negative securities positions are retained in an exception register; they receive zero allocation weight/cap in the nonnegative representation while aggregate controls remain binding. A matrix zero means that no positive position is assigned within the retained coverage. It need not mean that the economic exposure is absent.

The reconstruction checks the long register against every matrix cell; loan and deposit positions against their source controls; group removal against gross bounds; bank balance-sheet identities; and all sovereign/security issuer controls. The government columns equal EUR 2,133.0bn (Germany), 2,954.1bn (France), 2,486.9bn (Italy), and 1,452.3bn (Spain). These totals cover debt securities, not all Maastricht debt. All reconciliation records and full-precision matrices accompany the replication files.

\FloatBarrier
